\documentclass{ctexart}
\usepackage{geometry}
\usepackage{caption}
\usepackage{multirow}
\usepackage{array}
\usepackage{amsmath}
\newcolumntype{C}[1]{>{\centering\arraybackslash}p{#1}}
\geometry{a4paper,left=1.5cm,right=1.5cm,top=2cm,bottom=2cm}
\captionsetup[table]{labelsep=space,justification=centering}
\begin{document}
\begin{table}[ht]
\centering
\caption*{表2 \hspace{3em}各峰所对应的栅极电压$V_{G2K}$}
\renewcommand\arraystretch{1.1}
\begin{tabular}{C{5em}|C{4em}|C{4em}|C{4em}|C{4em}|C{4em}}
\hline
\hline
峰 & 1 & 2 & 3 & 4 & 5 \\
\hline
$V_{G2K}$(V) &  &  &  &  &  \\
\hline
\hline
峰 & 6 & 7 & 8 & 9 & 10 \\
\hline
$V_{G2K}$(V) &  &  &  &  &  \\
\hline
\hline
\end{tabular}
\end{table}
\begin{table}[ht]
\centering
\caption*{表3 \hspace{3em}$V_0$的估计值$\Delta V_{G2K}$}
\renewcommand\arraystretch{1.1}
\begin{tabular}{C{5em}|C{4em}|C{4em}|C{4em}|C{4em}|C{4em}}
\hline
\hline
$\Delta V_{G2K}$ & 1 & 2 & 3 & 4 & 5 \\
\hline
$\hat{V_0}$(V) &  &  &  &  &  \\
\hline
\hline
$\Delta V_{G2K}$ & 6 & 7 & 8 & 9 &  \\
\cline{1-5}
$\hat{V_0}$(V) &  &  &  &  &  \\
\hline
\hline
\end{tabular}
\end{table}
\end{document}